\begin{document}

\begin{center}
    \fbox{試験開始日時(XXXX/XX/XX 20:00)まで，このPDFファイルの中を見てはいけません．}
\end{center}
\begin{center}
    {\huge \textbf{A　r　c　a　e　a　共　通　テ　ス　ト}}\\
    \vspace{1em}
    {\Large \textbf{主催者：tzug（@tzug\_c）}}
\end{center}
{\Large
\begin{enumerate}[I]
    \item \textbf{注　意　事　項}
    \begin{enumerate}[1]
        \item 出題科目，配点，大問数及び選択方法は下表の通りです．\\

        \begin{table}[h]
            \centering
            \begin{tabular}{|c|c|c|c|}
                \hline
                \multirow{2}{*}{　出　題　科　目　} & \multirow{2}{*}{　配　　点　} &\multirow{2}{*}{問　　題} & \multirow{2}{*}{　選　択　方　法　} \\
                &&& \\ \hline
                \multirow{2}{*}{A~r~c~a~e~a} &\multirow{2}{*}{30~点} & 　第　1　問　 & \multirow{2}{*}{必　　答}\\  
                    \cline{3-3}
                    &&　第　2　問　& \\
                    \hline
                    \multirow{2}{*}{文　　科} &\multirow{2}{*}{35~点} &　第　1　問　 & \multirow{2}{*}{必　　答}\\
                    \cline{3-3}
                    &&　第　2　問　& \\
                    \hline
                    \multirow{2}{*}{理　　科} &\multirow{2}{*}{35~点} & 　第　1　問　 & \multirow{2}{*}{必　　答}\\
                    \cline{3-3}
                    &&　第　2　問　& \\
                    \hline
            \end{tabular}
        \end{table}

        \item 試験中に，問題用紙の落丁・乱丁や，フォームの解答欄の欠損等に気付いた場合は，tzug（@tzug\_c）の~DM~まで知らせてください．
        \item 不正行為
        \begin{enumerate}[\ajMaru{\arabic*}]
            \item 解答にあたって，解答の作成を直接の目的とした生成AIの利用はしてはいけません．
            \item 解答に直接結びつく行為を目的としない範囲に限り，用語や背景事項の確認のための検索エンジンの使用は認めます．
        \end{enumerate}
        \item その他
        \begin{enumerate}[\ajMaru{\arabic*}]
            \item この注意事項は，試験終了後における本試験の感想や雰囲気等についての発信を妨げるものではありません．
            広く歓迎します．
        \end{enumerate}
    \end{enumerate}
    \vspace{2em}
    \textbf{解答上の注意は次のページに記載しています．必ず読みなさい．}
    \thispagestyle{empty}
    \clearpage
    \addtocounter{page}{-1}
    \newpage
    \item \textbf{解　答　上　の　注　意}
    \begin{enumerate}[1]
        \item 解答はすべて，定められたフォームに解答しなさい．
        \item 解答は，この問題用紙の問題番号に対応した解答欄に解答しなさい．
        \item 問題については，すべて~\textbf{Ver~X.XX.X~(XXXX/XX/XX)~時点}の情報に基づいて解答するものとし，それ以降の仕様変更等に依りません．
        \item 問題文中の~\textgt{\fgb{ア}}~，~\textgt{\fgb{イウ}}~などには，符号($-$，$\pm$)あるいは数字(0~～~9)が入り，
            \textgt{ア}，\textgt{イ}，\textgt{ウ}，…~の一つ一つはこれらのいずれかに対応します．
        \item 解答に際しては，必ず各問いにおいて指定された解答形式に従いなさい．
        例えば，答えが $2.5$ となる場合であっても，$\dfrac{\textgt{\fgb{カ}}}{\textgt{\fgb{オ}}}$~と指示されているときは，$\dfrac{5}{2}$~と解答しなさい．
        \item $\dfrac{10}{4}$~や~$2\sqrt{8}$~など，約分されていないものや，根号の中に現れる自然数が最小とならない形での解答は，正解と認められません．
        \item 分数型での解答を求められた場合，分数の符号は分子に付けなさい．
        \item 小数型での解答を求められた場合，指定された桁数の一つ下の桁を四捨五入して答えなさい．
        例えば，$\textgt{\fgb{キ}}~.~\textgt{\fgb{クケ}}$~に~$2.495$ と解答する場合は，$2.50$ と解答しなさい．
        \item 問題文中の二重四角で示された~\textgt{\dgb{ナ}}~などの欄には，選択肢の中から最も適切なものを一つ選んで答えなさい．
        \item 同一の問題文中に~\textgt{\fgb{ハヒ}}，~\textgt{\dgb{フ}}~などが2回以上現れる場合，2回目以降は~\textgt{\fb{ハヒ}}~，~\textgt{\db{フ}}~のように細字で記します．
    \end{enumerate}
\end{enumerate}
}
\thispagestyle{empty}
\clearpage
\addtocounter{page}{-1}
\end{document}